\textbf{Phần 1. Trắc nghiệm lựa chọn}
\setcounterpageref{ex}{0}

\begin{ex} %[4H6V31]
Hình bình hành có
\choice
{4 cạnh bằng nhau}
{4 góc vuông}
{2 đường chéo bằng nhau}
{\True 2 cặp cạnh đối diện song song và bằng nhau}
\loigiai{
Hình bình hành có hai cặp cạnh đối diện song song và bằng nhau.\\
\textbf{Đáp án của câu hỏi là D}
}
\end{ex}

\begin{ex} %[4H6V31]
Hình thoi là hình tứ giác có
\choice
{4 góc vuông}
{\True 4 cạnh bằng nhau và 2 cặp cạnh đối diện song song}
{2 đường chéo bằng nhau}
{Không có cạnh nào song song}
\loigiai{
Hình thoi có 4 cạnh bằng nhau và các cặp cạnh đối diện song song.\\
\textbf{Đáp án của câu hỏi là B}
}
\end{ex}

\begin{ex} %[4H6V31]
Hai đường chéo của hình thoi có đặc điểm là
\choice
{Bằng nhau}
{Song song với nhau}
{\True Vuông góc với nhau và cắt nhau tại trung điểm của mỗi đường}
{Không cắt nhau}
\loigiai{
Hai đường chéo của hình thoi vuông góc với nhau tại trung điểm mỗi đường.\\
\textbf{Đáp án của câu hỏi là C}
}
\end{ex}

\begin{ex} %[4H6V31]
Chu vi hình thoi có độ dài cạnh là $a$ được tính theo công thức
\choice
{\True $P = a \times 4$}
{$P = a + 4$}
{$P = a \times a$}
{$P = a \times 2$}
\loigiai{
$P = a \times 4$.\\
\textbf{Đáp án của câu hỏi là A}
}
\end{ex}

\begin{ex} %[4H6V31]
Một hình thoi có cạnh dài $8\text{ cm}$. Chu vi của hình thoi đó là
\choice
{$16\text{ cm}$}
{$64\text{ cm}$}
{$24\text{ cm}$}
{\True $32\text{ cm}$}
\loigiai{
$P = 8 \times 4 = 32\text{ cm}$.\\
\textbf{Đáp án của câu hỏi là D}
}
\end{ex}

\begin{ex} %[4H6V31]
Một hình bình hành có độ dài đáy là $a$, chiều cao là $h$ (cùng đơn vị đo). Diện tích hình bình hành là
\choice
{$S = (a + h) \times 2$}
{\True $S = a \times h$}
{$S = a \times h : 2$}
{$S = a + h$}
\loigiai{
$S = a \times h$.\\
\textbf{Đáp án của câu hỏi là B}
}
\end{ex}

\begin{ex} %[4H6V31]
Một hình bình hành có độ dài đáy $15\text{ cm}$, chiều cao tương ứng $8\text{ cm}$. Diện tích hình bình hành đó là
\choice
{$60\text{ cm}^2$}
{$46\text{ cm}^2$}
{\True $120\text{ cm}^2$}
{$23\text{ cm}^2$}
\loigiai{
$S = 15 \times 8 = 120\text{ cm}^2$.\\
\textbf{Đáp án của câu hỏi là C}
}
\end{ex}

\begin{ex} %[4H6V31]
Một mảnh đất hình bình hành có đáy dài $24\text{ m}$, chiều cao bằng một nửa độ dài đáy. Diện tích mảnh đất đó là
\choice
{\True $288\text{ m}^2$}
{$144\text{ m}^2$}
{$576\text{ m}^2$}
{$72\text{ m}^2$}
\loigiai{
Chiều cao: $24 : 2 = 12\text{ m}$.\\
Diện tích: $24 \times 12 = 288\text{ m}^2$.\\
\textbf{Đáp án của câu hỏi là A}
}
\end{ex}

\begin{ex} %[4H6V31]
Một hình thoi có chu vi bằng $60\text{ cm}$. Độ dài cạnh của hình thoi là
\choice
{$12\text{ cm}$}
{$20\text{ cm}$}
{$30\text{ cm}$}
{\True $15\text{ cm}$}
\loigiai{
Cạnh: $60 : 4 = 15\text{ cm}$.\\
\textbf{Đáp án của câu hỏi là D}
}
\end{ex}

\begin{ex} %[4H6V31]
Hình vuông có phải là một hình thoi đặc biệt không?
\choice
{Không phải}
{\True Phải, vì hình vuông có 4 cạnh bằng nhau}
{Chỉ đúng khi có 2 đường chéo bằng nhau}
{Không xác định được}
\loigiai{
Hình vuông là hình thoi đặc biệt có 4 góc vuông.\\
\textbf{Đáp án của câu hỏi là B}
}
\end{ex}

\begin{ex} %[4H6V31]
Một hình bình hành có chu vi $48\text{ cm}$, độ dài một cạnh là $14\text{ cm}$. Độ dài cạnh kề với nó là
\choice
{$8\text{ cm}$}
{$12\text{ cm}$}
{\True $10\text{ cm}$}
{$20\text{ cm}$}
\loigiai{
Nửa chu vi: $48 : 2 = 24\text{ cm}$.\\
Cạnh kề: $24 - 14 = 10\text{ cm}$.\\
\textbf{Đáp án của câu hỏi là C}
}
\end{ex}

\begin{ex} %[4H6V31]
Một mảnh vườn hình bình hành có đáy $30\text{ m}$, chiều cao $15\text{ m}$. Người ta trồng bắp cải trên mảnh vườn đó, cứ $1\text{ m}^2$ thu hoạch được $4\text{ kg}$ bắp cải. Cả mảnh vườn thu hoạch được số tạ bắp cải là
\choice
{\True $18\text{ tạ}$}
{$16\text{ tạ}$}
{$20\text{ tạ}$}
{$180\text{ tạ}$}
\loigiai{
Diện tích: $30 \times 15 = 450\text{ m}^2$.\\
Số bắp cải: $450 \times 4 = 1800\text{ kg} = 18\text{ tạ}$.\\
\textbf{Đáp án của câu hỏi là A}
}
\end{ex}

\textbf{Phần 2. Trắc nghiệm đúng sai}
\setcounterpageref{ex}{0}

\begin{ex} %[4H6V31]
Xét tính đúng sai về hình bình hành và hình thoi:
\choiceTF
{\True Hình bình hành có 2 cặp cạnh đối diện song song và bằng nhau.}
{\True Hình thoi có 4 cạnh bằng nhau.}
{Hai đường chéo của hình bình hành luôn vuông góc với nhau.}
{\True Hai đường chéo của hình thoi vuông góc với nhau.}
\loigiai{
\begin{itemize}
    \item Hình bình hành có 2 cặp cạnh đối song song và bằng nhau: Đúng.
    \item Hình thoi có 4 cạnh bằng nhau: Đúng.
    \item Đường chéo hình bình hành nói chung không vuông góc (chỉ hình thoi/vuông mới vuông góc). Sai.
    \item Hai đường chéo hình thoi vuông góc: Đúng.
\end{itemize}
\begin{itemize}
    \item a) \textbf{Đúng}.
    \item b) \textbf{Đúng}.
    \item c) \textbf{Sai}.
    \item d) \textbf{Đúng}.
\end{itemize}
}
\end{ex}

\begin{ex} %[4H6V31]
Một hình thoi có cạnh $12\text{ cm}$ và một hình vuông có cạnh $12\text{ cm}$. Xét tính đúng sai:
\choiceTF
{\True Chu vi hình thoi bằng chu vi hình vuông ($48\text{ cm}$).}
{\True Cả hai hình đều có 4 cạnh bằng nhau.}
{Hình thoi luôn có 4 góc vuông giống hình vuông.}
{\True Hai đường chéo của cả hai hình đều vuông góc với nhau.}
\loigiai{
\begin{itemize}
    \item Cùng chu vi $12 \times 4 = 48\text{ cm}$: Đúng.
    \item 4 cạnh bằng nhau: Đúng.
    \item Hình thoi không nhất thiết có 4 góc vuông. Sai.
    \item Hai đường chéo đều vuông góc: Đúng.
\end{itemize}
\begin{itemize}
    \item a) \textbf{Đúng}.
    \item b) \textbf{Đúng}.
    \item c) \textbf{Sai}.
    \item d) \textbf{Đúng}.
\end{itemize}
}
\end{ex}

\textbf{Phần 3. Trả lời ngắn}
\setcounterpageref{ex}{0}

\begin{ex} %[4H6V31]
Tính chu vi của hình thoi có cạnh dài $14\text{ cm}$.
\shortans{56}
\loigiai{
\textbf{Đáp số:} $56$\\
$14 \times 4 = 56\text{ cm}$.
}
\end{ex}

\begin{ex} %[4H6V31]
Một hình bình hành có đáy $18\text{ cm}$, chiều cao $10\text{ cm}$. Tính diện tích theo xăng-ti-mét vuông.
\shortans{180}
\loigiai{
\textbf{Đáp số:} $180$\\
$18 \times 10 = 180\text{ cm}^2$.
}
\end{ex}

\begin{ex} %[4H6V31]
Một hình thoi có chu vi bằng $84\text{ dm}$. Tính độ dài một cạnh theo đề-xi-mét.
\shortans{21}
\loigiai{
\textbf{Đáp số:} $21$\\
$84 : 4 = 21\text{ dm}$.
}
\end{ex}

\begin{ex} %[4H6V31]
Một hình bình hành có diện tích $200\text{ cm}^2$, chiều cao $8\text{ cm}$. Tính độ dài đáy tương ứng theo xăng-ti-mét.
\shortans{25}
\loigiai{
\textbf{Đáp số:} $25$\\
$200 : 8 = 25\text{ cm}$.
}
\end{ex}

\begin{ex} %[4H6V31]
Một mảnh đất hình bình hành có hai cạnh kề dài $15\text{ m}$ và $25\text{ m}$. Tính chu vi mảnh đất theo mét.
\shortans{80}
\loigiai{
\textbf{Đáp số:} $80$\\
$(15 + 25) \times 2 = 80\text{ m}$.
}
\end{ex}

\begin{ex} %[4H6V31]
Một hình bình hành có độ dài đáy gấp đôi chiều cao và diện tích là $72\text{ cm}^2$. Tính chiều cao của hình bình hành theo xăng-ti-mét.
\shortans{6}
\loigiai{
\textbf{Đáp số:} $6$\\
$a = 2h \Rightarrow 2h \times h = 72 \Rightarrow h \times h = 36 \Rightarrow h = 6\text{ cm}$.
}
\end{ex}
